\lecture{9}{Tue 20 Sep 2022 12:03}{Integration}

A \textit{curve} in region $\Omega$ is a continuous map \[
	\gamma: [a,b] \to  \Omega \subset \mathbb{C}
.\] 

We can write it in term of real and imaginary parts, 
\[
	\gamma(t) = x(t) + i y(t)
.\] 

We consider \textit{piecewise-smooth} curves, curves such that there exists \[
	a < = t_0 < t_1 < \ldots < t_{n+1} = b
\] such that $\gamma$ has a continuous analytic on each interval $[t_i, t_{i+1}]$.

A curve that does not intersect itself is called \textit{simple}, i.e., $\gamma$ is injective.

A curve such that $\gamma(a) = \gamma(b)$ is called \textit{closed}.

If $\gamma$ is closed and injective on $[a,b)$, then the curve is called \textit{Jordan}.

(This theorem is not relevant to this class.)
\begin{theorem}[Jordan curve theorem]
	A Jordan curve divides a plane to two parts, inside and outside.

	Any simple closed contour separates the plane into two domains, each having the curve as its boundary. One of these domains, called the interior, is bounded; the other, called the exterior, is unbounded.
\end{theorem}

\begin{definition}[Line Integral]
\[
	\int_\gamma f(z) dz := \int_a^b f(\gamma(t))\gamma'(t)dt
.\] 
If $f(z) = u(x,y) i v(x,y)$, then $\gamma = x(t) + i y(t)$. Then the integral can be described in real integral:
\[
	\int_a^b (u + i v)(dx + i dy) = \int_\gamma u dx - v dy + i\int_\gamma udy + v dx 
.\] 
\end{definition}

\begin{example}
	Compute \[
		\int_\gamma |z| dz, \quad \gamma (t) = (1+i)t, \quad t \in [0,1]
	.\] 

	\begin{align*}
		&= \int _0^1 \sqrt{2 t^2}(1+i)dt  \\
		&= \sqrt{2} (1+i) \int _0^1 t dt \\
		&= \frac{1+i}{\sqrt{2} }
	.\end{align*}
\end{example}

\begin{proposition}(Properties of integral)
	\begin{enumerate}
		\item Suppose we have 
			\[
				\varphi: [c,d] \to  [a,b]
			.\] 
			And \[
				\gamma: [a,b] \to  \Omega
			.\] 
		Where $\varphi(c) = a, \varphi(d) = b$, $\varphi$ is strictly increasing. Then
		\[
			\int_{\varphi\circ\gamma}f(z) dz = \int_\gamma f(z) dz
		.\] 
	\item (Linearity)
		\[
		\int _\gamma (c_1f + c_2g)dz 
		=c_1 \int_\gamma fdz + c_2 \int _\gamma g dz 
		.\] 
	\item (Addition of paths)
		If $\gamma = \gamma_1 + \gamma_2$, then
		\[
		\int _\gamma f(z) dz
		= \int _{\gamma_1} f(z) dz
		+ \int _{\gamma_2} f(z) dz
		.\] 
	\item 
	\end{enumerate}
\end{proposition}

\begin{definition}
	Assume $f$ is continuous in $\Omega$. If there exists an analytic $F$ in $\Omega$ such that $F' = f$, we say that $f$ has a \textit{primitive} (indefinite integral) in $\Omega$.
\end{definition}

Observe that we have by Newton Leibnitz formula: \[
	\int _\gamma f(z) dz = \int _\gamma F'(z) dz = \int _a^b F'(\gamma(t))\gamma'(t) dt = F(\gamma(b)) - F(\gamma(a))
.\] 

Hence the integral is \textit{path independent} if the function has a primitive.

\begin{example}
	$\Omega = \mathbb{C}^*$, and $f(z) = z^n$. Consider any integer $n$.

	For $n \neq  -1$ has a primitive $F(z) = \frac{1}{n+1}z^{n+1}$. So any integral is path independent.	
\end{example}

\begin{theorem}
	For a continuous function $f$ in a region $\Omega$, the following are equivalent:
	\begin{enumerate}
		\item $f$ has a primitive
		\item $\int _\gamma f dz$ is path independent for paths in $\Omega$
		\item $\int _\gamma f dz$ for every closed path is zero.
	\end{enumerate}
\end{theorem}

If $\int f(z) dz $ is path independent, we can define \[
	F(z) = \int _{z_0}^z f(\zeta)d\zeta
.\] 

Then we show by definition that $F$ has a derivative 
\begin{align*}
	F'(z) &= \lim_{h \to 0} \frac{F(z+h) - F(z)}{h}\\
	&= \lim_{h \to 0} \frac{\int _z^{z+h}f(z) dz }{h}\to f(z)
.\end{align*} 


Main estimate of an integral:
\[
	\left| \int _\gamma f(z) dz  \right| \le \max_{z \in \gamma} \left| f(z) \right| \cdot \text{length of }\gamma
.\] 

\begin{example}(The most important example)
	Consider
	\[
		\int _\gamma \frac{dz }{z} \quad, \gamma \text{ is unit circle, parameterized counter-clockwise}
	.\] 
	\[
		\gamma(t) = e^{it },\quad 0\le t\le 2\pi
	.\] 
	
	By definition,
	\begin{align*}
		&=\int _0^{2 \pi} \frac{i e^{it }}{e^{it }} dt\\
		&= 2 \pi i
	.\end{align*} 

	$f(z) = \frac{1}{z}$ is analytic in $\mathbb{C}^* = \Omega$. 
	$\gamma$ is a closed curve around $0$. However the integral is not zero.

	Hence $f$ \textbf{does not have a primitive} in $\Omega$.

	Consider $F(z) = \operatorname{Log} z$, $F'(z) = \frac{1}{z}$ is defined in $\mathbb{C} \setminus  (-\infty, 0] = \Omega_1$. 
	So $f(z) = \frac{1}{z}$ does have a primitive in $\Omega_1$. So whether having a primitive is a property depending on not only the function but also the region we choose.
\end{example}

\subsection{Calculus: Green's Formula}

Consider $\mathcal{D} = \{(x,y): a \le x\le b, \varphi_1(x)\le y\le \varphi_2(x)\} $.

Suppose we have some function $u(x,y)$ continuously differentiable in $\overline{\mathcal{D}} $. We apply Newton Leibneitz,
\[
	\int_{\varphi_1(x)}^{\varphi_2(x)}u_y dy = u(\varphi_2(x))-u(\varphi_1(x)) \text{ for } a\le x\le b
.\] 

Hence \[
	\iint_{\mathcal{D}} u_y dx dy = \int  u(\varphi_2(x))dx-\int u(\varphi_1(x)) dx
.\] 

Assume $\partial \mathcal{D}$ is a parameterized curve (oriented boundary): when described, the region is on the left.

\[
	= -\int_{\partial \mathcal{D}} u(x,y) dx
.\] 


Similarly for $u_x$ ($\varphi_1$, $\varphi_2$, $a$ and $b$ are re-defined):
\[
	\int_a^b\int _{\varphi_1(y)}^{\varphi_2(y)} u_x dx dy
	= \int_{\partial \mathcal{D}}u dy
.\] 

\begin{definition}[Green's Formula]
\[
	\int _{\partial \mathcal{D}} u(x,y)dx + v(x,y)dy
	= \iint_{\mathcal{D}} - u_y + v_x dA
.\] 
\end{definition}

Suppose we have a region split by some curve $\gamma_3$, as shown in the following figure: FIGURE
\[
	\iint_{\mathcal{D}}=
	\iint_{\mathcal{D}_1}+
	\iint_{\mathcal{D}_2}
.\] 
Our definition of oriented boundary gives the next property for line integrals:
\[
	\int _{\partial \mathcal{D}}=
	\int _{\partial \mathcal{D}_1}+
	\int _{\partial \mathcal{D}_2}
.\] 
Since
\[
	\partial \mathcal{D}_1
	+\partial \mathcal{D}_2
	= -\gamma_2 + \gamma_3 + (\gamma_1-\gamma_3) = \gamma_1 - \gamma_2 = \partial \mathcal{D}
.\] 


Now we can apply this argument of cutting to very complex regions by cutting it into simple regions.

If C-R is satisfied, then any integral around closed loop is zero. Such differential forms are called \textit{exact}. This is the first condition of path independence.

If $u dx + v dy $ is a exact form in $\Omega$, then 
$\int _{\partial \mathcal{D}} u dx + v dy = 0$ for $\mathcal{D} \subset \Omega$.
